\section{Polynomial-weighted causal synthesis for actual character PW levels}
\label{sec:character-weighted-pw}

The nontrivial character source has no archimedean polar moment
constraint.  Its reflected theta formula nevertheless has an exact
conductor scale and a dual character.  We retain both, identify the
actual weighted range by a bounded causal convolution, and prove
the spectral and ambient trace assertions in that quotient.
No weighted Beurling identification is assumed.

Let \(K=\mathbb Q\) or \(\mathbb Q(\zeta_3)\), and let \(\eta\)
be a nonprincipal primitive Dirichlet character in the first case,
or a nontrivial finite-order Hecke character covered by
MA-CGP-001 in the second.  This includes \(\psi_{269,3}\).
Write
\[
 \begin{gathered}
 D_\infty=[K:\mathbb Q],\quad
 w_K=|\mathcal O_K^\times|,\quad Q=|D_K|N\mathfrak f_\eta,\\
 d=\tfrac12-\alpha,\quad u=\tfrac12+\alpha,
 \qquad 0<\alpha<\tfrac12,\quad s\ge0 .
 \end{gathered}
 \tag{MA-GW.45}
\]
Here \(D_\infty=1,2\) and \(w_K=2,6\), respectively.
The degree \(D_\infty\) is distinct from the strip parameter \(d\).
For the primitive archimedean norm band \(\Lambda>0\), set
\[
 Y=Q\Lambda,\quad c=-\log Y,\quad
 H=L^2(\mathbb R,e^{2\alpha|x|}(1+x^2)^{s/2}dx),\quad
 H_c=\mathbf1_{[c,\infty)}H .
 \tag{MA-GW.46}
\]
Let \(J_{\eta,\Lambda}\) be the actual Hilbert closure of
\(\mathscr P_{\eta,\Lambda}\) in (MA-CGP.9), written in the
logarithmic variable.  Its source invariant level is exactly
\(Y\), with \(S'=S\), by (MA-CGP.11); no extra unit-shell
Euler deletion is inserted.  Let \(Q_{\eta,\Lambda}\) be its
orthogonal projection in \(H\), and put
\[
 \mathcal Q_{\eta,c}=H_c/J_{\eta,\Lambda},\qquad
 Z_{\eta,c}=H_c\ominus J_{\eta,\Lambda},\qquad
 P_{\eta,c}^Z=\text{the orthogonal projection onto }Z_{\eta,c}.
 \tag{MA-GW.47}
\]
The last projection is extended by zero on \(H_c^\perp\).

\paragraph{MA-GW-004 (the actual character theorem at every polynomial weight).}
Let \(\rho\) run through the zeros of the primitive entire
completion \(\Lambda_\eta\), with multiplicities \(m_\eta(\rho)\).
Define
\[
 b_s=\#\{j\in\mathbb Z_{\ge0}:2j<s-1\},\qquad
 r^+_{\eta,\rho}=
 \begin{cases}
 m_\eta(\rho),&\operatorname{Re}\rho<u,\\
 \min(m_\eta(\rho),b_s),&\operatorname{Re}\rho=u,\\
 0,&\operatorname{Re}\rho>u .
 \end{cases}
 \tag{MA-GW.48}
\]
The actual closed subspace \(J_{\eta,\Lambda}\subset H_c\)
is precisely the common kernel of the continuous Mellin jets
of orders \(0\le j<r^+_{\eta,\rho}\) at \(\rho-\tfrac12\).
The corresponding finite spectral blocks span a dense subspace
of the actual quotient.

For every additive \(g\in C_c^\infty((0,\infty))\), both
actual operators
\[
 C_g^{\mathcal Q_{\eta,c}}\in\mathcal S_1(\mathcal Q_{\eta,c}),
 \qquad
 C_gP_{\eta,c}^Z\in\mathcal S_1(H)
 \tag{MA-GW.49}
\]
are trace class, and
\[
 \operatorname{Tr}_{\mathcal Q_{\eta,c}}C_g^{\mathcal Q_{\eta,c}}
   =\operatorname{Tr}_H(C_gP_{\eta,c}^Z)
   =\sum_\rho r^+_{\eta,\rho}\,\mathcal Mg(\rho-\tfrac12),
 \qquad
 \mathcal Mg(z)=\int_0^\infty g(t)e^{zt}dt .
 \tag{MA-GW.50}
\]
The sum is absolutely convergent and independent of \(\Lambda\).
For the uniformly bounded smooth compact exhaustion
\(\chi_R(x)=\chi_0(x/R)\), with \(\chi_0=1\) near zero,
\[
 \begin{aligned}
 \lim_{\Lambda\to\infty}\lim_{R\to\infty}
    \operatorname{Tr}(M_{\chi_R}C_gQ_{\eta,\Lambda})
       &=-\sum_\rho r^+_{\eta,\rho}\mathcal Mg(\rho-\tfrac12),\\
 \lim_{R\to\infty}\lim_{\Lambda\to\infty}
    \operatorname{Tr}(M_{\chi_R}C_gQ_{\eta,\Lambda})
       &=-\sum_\rho r^{\rm strip}_{\eta,\rho}
                                      \mathcal Mg(\rho-\tfrac12),
 \end{aligned}
 \tag{MA-GW.51}
\]
where
\[
 r^{\rm strip}_{\eta,\rho}=
 \begin{cases}
 m_\eta(\rho),&d<\operatorname{Re}\rho<u,\\
 \min(m_\eta(\rho),b_s),&
            \operatorname{Re}\rho=d\text{ or }\operatorname{Re}\rho=u,\\
 0,&\text{otherwise}.
 \end{cases}
 \tag{MA-GW.52}
\]
The first limit minus the second equals
\[
 -\sum_{\operatorname{Re}\rho<d}
             m_\eta(\rho)\mathcal Mg(\rho-\tfrac12)
 -\sum_{\operatorname{Re}\rho=d}
        \bigl(m_\eta(\rho)-\min(m_\eta(\rho),b_s)\bigr)
                                  \mathcal Mg(\rho-\tfrac12).
 \tag{MA-GW.53}
\]
The theorem asserts neither the existence of such omitted zeros
or jets nor the nonvanishing of this expression for a given test.
It uses no \(\mathsf A\), \(M\), generic-edge assumption, or
separation between zeros.

\begin{proof}
\emph{1. The actual source and its bounded causal symbol.}
The map
\[
 (U_cf)(y)=e^{\alpha(c+y)}f(c+y),\qquad
 W_c(y)=e^{2\alpha|c+y|-2\alpha(c+y)}
                                      (1+(c+y)^2)^{s/2}
 \tag{MA-GW.54}
\]
is unitary from \(H_c\) onto \(H_{W_c}=L^2((0,\infty),W_cdy)\).
For this fixed \(c\), that norm is equivalent to
\(\mathcal H_s=L^2((0,\infty),(1+y^2)^{s/2}dy)\).
The polynomial weight has not been removed.

Choose an integer \(N\) sufficiently large and set
\[
              \Phi_\eta(z)=
                 \frac{L_K(z+d,\bar\eta)}{(z+2)^N}.
 \tag{MA-GW.55}
\]
There is no pole-canceling numerator factor: the nontrivial
finite Hecke function is entire.  Put \(\epsilon_0=d/2\).
The fixed-character vertical polynomial bound of NT-Z10,
and the absolutely convergent Euler product on a right
half-plane, permit \(N>3\) with
\[
 \begin{aligned}
 |\Phi_\eta(\sigma+it)|&\le C(1+|t|)^{-2}
                  &&(-\epsilon_0\le\sigma\le2),\\
 |\Phi_\eta(z)|&\le C|z|^{-N}
                  &&(\operatorname{Re}z\ge2).
 \end{aligned}
 \tag{MA-GW.56}
\]
There is no singularity on this closed half-plane.
The inverse Laplace integral may be moved to
\(\operatorname{Re}z=-\epsilon_0\) by finite rectangles
and the uniform vertical decay.  It gives a continuous
kernel \(\phi_\eta\) with
\(|\phi_\eta(t)|\le Ce^{-\epsilon_0t}\) for \(t\ge0\).
For \(t<0\), move instead to \(\operatorname{Re}z=M\).
The absolute integral is at most \(Ce^{Mt}M^{1-N}\);
letting \(M\to\infty\) proves that the kernel is causal.
Fourier inversion on a vertical line, followed by analytic
continuation, gives \(\mathcal L\phi_\eta=\Phi_\eta\) on
\(\operatorname{Re}z>-\epsilon_0\).

The unilateral shifts \(S_tf(y)=\mathbf1_{y\ge t}f(y-t)\)
on \(\mathcal H_s\) are strongly continuous and satisfy
\[
 \|S_t\|\le2^{s/4}(1+t^2)^{s/4},\qquad
 T_{\phi_\eta}f=\int_0^\infty\phi_\eta(t)S_tf\,dt
                                      \in\mathcal B(\mathcal H_s).
 \tag{MA-GW.57}
\]
The inequality follows from
\(1+(y+t)^2\le2(1+y^2)(1+t^2)\).
Continuity follows on compact smooth functions and then
by density.  The integral is vectorwise, with an integrable
operator-norm majorant.  Its Laplace action is multiplication
by \(\Phi_\eta\).

We now prove the exact range equality
\[
 U_cJ_{\eta,\Lambda}
        =\overline{\operatorname{ran}T_{\phi_\eta}}^{\,\mathcal H_s}.
 \tag{MA-GW.58}
\]
Let \(u_\infty=\widehat f\) be any actual smooth frequency
test of the required parity or radial type supported in
\(\{|\xi|_\infty\le\Lambda\}\), and put
\(U(r)=u_\infty(r^{1/D_\infty})\).
No value or moment condition at zero is imposed.
Indeed the nontrivial primitive finite vector in
(MA-CGP.5) already kills both adelic polar terms.
At least one finite place is ramified in these
class-number-one character families, as proved in
MA-CGP-001.

The actual reflected theta calculation (MA-CGP.13) gives
\[
 \begin{aligned}
 (\mathcal L U_c\Theta_\eta f)(z)
  &=w_K\epsilon_\eta^\theta e^{zc}Q^{z-\alpha}
       L_K(z+d,\bar\eta)\int_0^\Lambda U(r)r^{z+d-1}dr\\
  &=C_\eta(Y)\,L_K(z+d,\bar\eta)\,\mathcal La(z),\\
 a(q)&=e^{-dq}U(\Lambda e^{-q}),\qquad
 C_\eta(Y)=w_K\epsilon_\eta^\theta\,Y^dQ^{-1/2}.
 \end{aligned}
 \tag{MA-GW.59}
\]
For completeness the scalar simplification is
\(e^{zc}Q^{z-\alpha}\Lambda^{z+d}
=Q^{-\alpha}\Lambda^d=Y^dQ^{-1/2}\).
The first equality is also valid for a frequency test
nonzero at the origin: initially take \(\operatorname{Re}(z+d)>1\),
where the ideal sum and the integral converge absolutely.
Both sides then continue to \(\operatorname{Re}z>0\).
Thus no annular-density assertion is needed for this inclusion.

At \(q=0\), \(a\) and all its derivatives vanish because
the smooth compact frequency test is flat at the band edge.
At infinity its derivatives are \(O(e^{-dq})\).
For real parity this follows from the smooth positive-ray
function \(U\) and its logarithmic derivatives; for complex
radial smooth functions, \(u_\infty(\sqrt r)\) is smooth at
\(r=0\), and the same bound holds.  In particular it holds
even when \(U(0)\ne0\).
Consequently \((\partial+2)^Na\in\mathcal H_s\), all
integration-by-parts boundary terms at zero vanish, and
(MA-GW.59) becomes
\[
 U_c\Theta_\eta f
           =C_\eta(Y)\,T_{\phi_\eta}(\partial+2)^Na.
 \tag{MA-GW.60}
\]
Laplace uniqueness justifies this equality of actual
half-line functions.  This proves inclusion of every
actual source core vector in the bounded convolution range.

Conversely, any \(a\in C_c^\infty((0,\infty))\) gives an
actual frequency test by
\(U(\Lambda e^{-q})=e^{dq}a(q)\).
Its parity or radial extension is smooth, supported away
from the origin and the band edge, and imposes no moment
condition.  Its Fourier inverse belongs to the actual
archimedean Schwartz source.  Formula (MA-GW.60) applies.
For \(h\in C_c^\infty((0,\infty))\), solve
\((\partial+2)^Na=h\) with zero initial traces using the
causal kernel \(e^{-2q}q^{N-1}/(N-1)!\).
The solution is flat at zero and all its derivatives
decrease exponentially, up to a fixed polynomial factor.
Smooth compact truncations \(a_M\) satisfy
\((\partial+2)^Na_M\to h\) in \(\mathcal H_s\).
Boundedness of \(T_{\phi_\eta}\), the nonzero scalar
\(C_\eta(Y)\), and density of these \(h\) prove the reverse
inclusion in (MA-GW.58).  Equivalent fixed-\(c\) norms
preserve this same closed subspace.

\emph{2. The actual quotient resolvent and its character index.}
Set \(\mathcal J=\overline{\operatorname{ran}T_{\phi_\eta}}\)
and \(\mathcal Q=\mathcal H_s/\mathcal J\).
The shifts commute with \(T_{\phi_\eta}\), hence induce a
strongly continuous polynomially bounded semigroup on
\(\mathcal Q\).  Let \(A\) be its generator.
For \(\operatorname{Re}\lambda>0\) its actual resolvent is
\[
 R(\lambda)=\int_0^\infty e^{-\lambda t}S_t^{\mathcal Q}dt.
 \tag{MA-GW.61}
\]
The generator difference quotient, integrated against this
exponentially damped kernel, verifies both inverse identities
and the generator-domain inclusion.

For \(\operatorname{Re}\lambda<\epsilon_0\), define
\[
 \psi_\lambda(r)=\int_r^\infty
                  \phi_\eta(t)e^{\lambda(t-r)}dt,\qquad
 T_\lambda=\int_0^\infty\psi_\lambda(r)S_r^{\mathcal Q}dr.
 \tag{MA-GW.62}
\]
Differentiation under the exponentially dominated integral
makes this operator-norm holomorphic there.  Moreover
\[
 \sup_{\operatorname{Re}\lambda\le\epsilon}\|T_\lambda\|<\infty
                                      \qquad(\epsilon<\epsilon_0).
 \tag{MA-GW.63}
\]
An explicit bound is a constant times
\(\int_0^\infty|\phi_\eta(t)|e^{\max(\epsilon,0)t}
\int_0^t(1+r)^{s/2}dr\,dt\).
It is uniform in height and throughout that entire left
half-plane.

The identities
\(\psi_\lambda'=-\phi_\eta-\lambda\psi_\lambda\) and
\(\psi_\lambda(0)=\Phi_\eta(-\lambda)\) give
\[
 (\lambda-A)T_\lambda=\Phi_\eta(-\lambda)1,\qquad
 T_\lambda(\lambda-A)=\Phi_\eta(-\lambda)1
                                      \quad\text{on }\mathcal D(A).
 \tag{MA-GW.64}
\]
Before quotienting the right side is
\(\Phi_\eta(-\lambda)1-T_{\phi_\eta}\).
The first identity follows by a shift difference quotient
and integration by parts in \(r\); the weighted exponential
majorants above prove that the result lies in the generator
domain.  Applying the same calculation to a domain vector
proves the second identity.  All operators preserve
\(\mathcal J\), because their convolutions commute.
Thus, wherever the scalar is nonzero,
\[
                         R(\lambda)=T_\lambda/\Phi_\eta(-\lambda)
 \tag{MA-GW.65}
\]
is the actual two-sided inverse of the closed generator.
It agrees with (MA-GW.61) on their intersection.
These formulas give a meromorphic resolvent on the plane;
all apparent poles in the right half-plane are removable.

For \(\operatorname{Re}\lambda\le0\), the finite-function
argument \(d-\lambda\) has positive real part.
Its archimedean factors are nonzero and holomorphic:
they are \(\Gamma_{\mathbb R}(d-\lambda+e)\) in the real
case and \(\Gamma_{\mathbb C}(d-\lambda)\) in the complex
case, up to the nonzero conductor factor.
The entire completed functional equation
\(\Lambda_{\bar\eta}(w)=W_{\bar\eta}\Lambda_\eta(1-w)\)
therefore identifies every possible pole, including its
multiplicity, as
\[
 \lambda_\rho=\rho-u,\qquad \operatorname{Re}\rho\le u,
 \qquad\text{pole order at most }m_\eta(\rho).
 \tag{MA-GW.66}
\]
The divisor is that of \(\eta\), not \(\bar\eta\).
The classical containing strip places these possible
poles in \([-u,0]+i\mathbb R\).
No pole or artificial zero occurs at \(d-\lambda=1\);
the finite character function here is nontrivial and
entire.  Any zero there would be an actual completed
zero paired with the completed endpoint at zero.
Trivial finite-\(L\) zeros have \(d-\lambda\in\mathbb R_{\le0}\) in these
finite-order families and lie outside the left-half-plane
calculation.  No hidden gamma divisor is added.

\emph{3. Finite rank and the exact edge jets.}
At a possible pole \(\lambda_0\) of scalar multiplicity \(m\),
define \(E=(2\pi i)^{-1}\int_\gamma R(\lambda)d\lambda\)
on a small positive circle.  Integrating the resolvent
identity over nested circles gives \(E^2=E\) and commutation
with the resolvent.  Since \(AR(z)=zR(z)-1\), the contour
integrals also converge in each required graph norm and give
\[
 (A-\lambda_0)^mE=0,\qquad
 [(\lambda-\lambda_0)^{-j}]R(\lambda)
                                =(A-\lambda_0)^{j-1}E .
 \tag{MA-GW.67}
\]
These identities alone do not assert finite rank; we prove it.

The continuous complex-linear dual of \(\mathcal H_s\)
is represented by densities \(v\) through
\(\ell_v(f)=\int_0^\infty f(y)v(y)dy\), with
\[
                \int_0^\infty|v(y)|^2(1+y^2)^{-s/2}dy<\infty .
 \tag{MA-GW.68}
\]
The quotient dual is the annihilator of \(\mathcal J\).
Its transpose translations act by \(v(y+t)\).
Thus every vector in \(\operatorname{ran}E'\) satisfies
\((\partial-\lambda_0)^mv=0\) distributionally on the
connected half-line.  Multiplication by \(e^{-\lambda_0y}\)
and successive integration show that it has the form
\(e^{\lambda_0y}p(y)\), with \(\deg p<m\).
If \(\operatorname{Re}\lambda_0<0\), every such polynomial
satisfies (MA-GW.68).  If \(\operatorname{Re}\lambda_0=0\),
its highest degree \(j\) satisfies it exactly when
\(2j<s-1\).  Equality gives logarithmic divergence.
This bounds the dimension of \(\operatorname{ran}E'\)
by \(m\) or \(\min(m,b_s)\).  The range of a bounded
projection and the range of its transpose have the same
finite dimension once either is finite: restricting the
transpose identifies it with the full continuous dual of
the closed range.  Hence this is an actual finite-rank bound.

It is attained.  For each allowed \(0\le j<m\), the
continuous functional
\[
                    \ell_j(f)=\int_0^\infty
                                   y^j e^{\lambda_0y}f(y)dy
 \tag{MA-GW.69}
\]
annihilates \(T_{\phi_\eta}f\).
After substituting the causal convolution and expanding
\((y+t)^j\), each coefficient is a derivative of
\(\Phi_\eta(-\lambda)\) at \(\lambda_0\) of order at most
\(j<m\), and therefore vanishes.  The exponential kernel
bound and (MA-GW.68) justify all integrals.
These independent functionals form a finite triangular
chain under translations.  The Laplace resolvent
(MA-GW.61) on that chain is its ordinary finite-dimensional
resolvent; the small contour integral fixes each member.
They belong to \(\operatorname{ran}E'\), giving the
matching lower bound and its full dual range.
Equation (MA-GW.67) now proves finite rank of every
principal part.

Pulling (MA-GW.69) back by \(U_c\) gives, up to a nonzero
constant,
\[
 f\longmapsto\int_c^\infty
             (x-c)^j e^{(\lambda_0+\alpha)x}f(x)dx .
 \tag{MA-GW.70}
\]
Here \(\lambda_0+\alpha=\rho-\tfrac12\).
The invertible triangular change from \((x-c)^j\) to
\(x^j\) gives the stated Mellin jets and the lengths
(MA-GW.48).

\emph{4. Uniform contour bounds and actual spectral synthesis.}
Fix \(0<\epsilon<\epsilon_0\) and \(C_0>u+2\).
The right-line resolvent at \(\operatorname{Re}\lambda=\epsilon\)
is uniformly bounded by (MA-GW.57) and (MA-GW.61).
On the whole region \(\operatorname{Re}\lambda\le-C_0\),
the Euler product gives
\[
 |L_K(d-\lambda,\bar\eta)|
   \ge\prod_{\mathfrak p}(1-(N\mathfrak p)^{-(d+C_0)})
   =\zeta_K(d+C_0)^{-1}>0 .
\]
Together with (MA-GW.63) and (MA-GW.65), this yields
\[
                   \|R(\lambda)\|\le C(1+|\lambda|)^N
                         \qquad(\operatorname{Re}\lambda\le-C_0).
 \tag{MA-GW.71}
\]
The constant is uniform as the left line moves.
The exponent is \(N\), rather than \(N-1\) in the
pole-canceled zeta case.

The fixed-character good-height proof MT-M18,
(M18.5)--(M18.8), applies to \(L_K(\,\cdot\,,\bar\eta)\).
It uses the unconditional local count from NT-Z10,
the Euler lower bound at \(2+iT\), the logarithmic-distance
averaging on a half-unit interval, and the finite
disk-Blaschke/Harnack argument.  In particular it gives a
polynomial lower bound for the finite \(L\)-function along
one fixed horizontal strip segment, independently of
separations of zeros inside the blocks.
Reflecting the chosen heights under \(w=d-\lambda\)
gives \(t_n=n+O(1)\), with gaps between \(1/2\) and \(3/2\).
For
\(\Gamma_n=\partial([-C_0,\epsilon]+i[t_n,t_{n+1}])\),
all low contours adjusted individually if necessary,
put \(E_n=(2\pi i)^{-1}\int_{\Gamma_n}R(\lambda)d\lambda\).
We obtain
\[
 \sup_{\Gamma_n}\|R(\lambda)\|\le C(2+|n|)^B,\qquad
 \operatorname{rank}E_n=O(\log(2+|n|)).
 \tag{MA-GW.72}
\]
The right side uses the actual right resolvent if the
scalar denominator has a removable zero there.
Every horizontal side avoids all relevant zeros.
The rank bound follows from the finite-rank calculation
and the completed local zero count.
Only this one central strip uses the good-height estimate;
the expanding left region uses the uniform bound
(MA-GW.71).

For \(k\in C_c^\infty((0,\infty))\) put
\(H_k(\lambda)=\int k(t)e^{\lambda t}dt\) and
\[
 V_n(k)=\frac1{2\pi i}\int_{\Gamma_n}
                             H_k(\lambda)R(\lambda)d\lambda .
 \tag{MA-GW.73}
\]
Their ranges lie in \(\operatorname{ran}E_n\) by the
resolvent contour identities.  Repeated integration by
parts in \(H_k\) and (MA-GW.72) give
\[
 \sum_n\|V_n(k)\|_1
  \le C_M\sum_n\log(2+|n|)(2+|n|)^{B-M}<\infty
                                             \qquad(M>B+2).
 \tag{MA-GW.74}
\]
Trace class is therefore established before taking residues
or scalar traces.

On the line \(\operatorname{Re}\lambda=\epsilon\), inserting
the actual Laplace resolvent and applying Fourier inversion
gives
\[
 C_k^{\mathcal Q}
   =\frac1{2\pi i}\int_{\epsilon-i\infty}^{\epsilon+i\infty}
                                  H_k(\lambda)R(\lambda)d\lambda .
 \tag{MA-GW.75}
\]
Fubini is justified by the exponentially damped polynomial
semigroup bound and rapid vertical decay of \(H_k\).
Finite rectangles with good horizontal heights move this
integral to \(-C_0\) and produce precisely the block sum;
the horizontal integrals vanish by (MA-GW.72).
There are no poles farther left.  Move the remaining
line to \(-M\), using (MA-GW.71), and then let \(M\to\infty\).
If \(\operatorname{supp}k\subset[a,b]\), \(a>0\), then
\[
 |H_k(-M+it)|
   \le C_{\ell,k}e^{-Ma}(1+M)^\ell(1+|t|)^{-\ell}.
\]
For \(\ell>N+2\), (MA-GW.71) makes the remaining
operator-norm integral tend to zero.
Consequently the actual operator identity is
\[
                          C_k^{\mathcal Q}=\sum_nV_n(k)
                                      \quad\text{in }\mathcal S_1 .
 \tag{MA-GW.76}
\]
Choose a positive-time smooth approximate identity supported
in \([\varepsilon,2\varepsilon]\).
Its quotient convolutions tend strongly to the identity.
Equation (MA-GW.76) proves density of the finite blocks.
If every retained jet vanishes on a quotient vector, the
full dual ranges determined above show that every \(E_n\)
kills it; then every \(C_k^{\mathcal Q}\) kills it, and
the approximate identity makes it zero.
This proves the exact closed-kernel assertion, with no
remaining unsynthesized quotient.

On each finite block the trace is its length times
\(H_k(\lambda_\rho)\); all higher nilpotent terms have
trace zero.  Set \(k(t)=e^{\alpha t}g(t)\).
Then \(H_k(\lambda_\rho)=\mathcal Mg(\rho-\tfrac12)\).
The trace-norm series and bounded similarity to the
actual quotient norm prove its trace assertion in
(MA-GW.50).  Absolute convergence also follows from
NT-Z10 and rapid vertical Mellin decay on the fixed strip.

We also need an actual trace-class resolvent power.
Take \(\nu>\epsilon\) and an integer
\(q>\max\{B+2,N+2\}\).  Repeat the same contours with
\(H(\lambda)=(\nu-\lambda)^{-q}\).
On the right line its actual Laplace kernel is
\(t^{q-1}e^{-\nu t}/(q-1)!\), so Fourier inversion gives
\((\nu-A)^{-q}\).
The block series converges in trace norm by (MA-GW.72).
The far-left remainder is at most
\(CM^{N+1-q}\) by (MA-GW.71); the horizontal integrals
vanish by the same chosen inequalities.
Thus
\[
                            (\nu-A)^{-q}\in\mathcal S_1(\mathcal Q).
 \tag{MA-GW.77}
\]

\emph{5. The actual ambient lift, with its endpoint and cusp.}
Return by the bounded similarity to \(H_{W_c}\), and
write \(W=W_c\).  The adjoint shifts and their generator are
\[
 (S_t^{*,W}v)(y)=\frac{W(y+t)}{W(y)}v(y+t),\qquad
 L_Wv=W^{-1}(Wv)' .
\]
Multiplication by \(W\) is unitary onto \(L^2(W^{-1}dy)\),
where this is the left-translation generator with no
boundary condition at zero.  The fundamental theorem of
calculus gives its difference-quotient domain; the converse
follows by testing against compact interior functions.
Induction gives
\[
 v\in\mathcal D(L_W^q)
 \quad\Longleftrightarrow\quad
 (Wv)^{(j)}\in L^2(W^{-1}dy)\quad(0\le j\le q).
 \tag{MA-GW.78}
\]
The required translation majorants follow from (MA-GW.57)
and fixed-\(c\) norm equivalence.

The actual orthogonal complement of the invariant ideal
is invariant under these adjoints.
The full resolvent restricts to it by its Laplace integral.
Adjunction of (MA-GW.77) therefore gives
\[
 \mathcal R_q=(\nu-L_W)^{-q}P_{U_cZ_{\eta,c}}
                   \in\mathcal S_1(H_W),
 \tag{MA-GW.79}
\]
and it maps boundedly into the \(q\)th graph domain.
This last assertion follows from
\(L_W(\nu-L_W)^{-1}=\nu(\nu-L_W)^{-1}-1\).

To make the ambient factorization explicit, put \(a=W^{-1}\),
\(b=Wv\), and \(k=e^{\alpha\,\cdot}g\).  For a domain vector,
\[
 C_k^+(\nu-L_W)^qv(x)
       =\int_0^\infty k(x-y)a(y)(\nu-\partial_y)^qb(y)dy .
 \tag{MA-GW.80}
\]
Split at \(L=-c\) if \(c<0\).
On each smooth piece every \(W\partial^j(W^{-1})\)
is bounded, because all positive-order derivatives of
\[
 2\alpha|c+y|-2\alpha(c+y)+\tfrac s2\log(1+(c+y)^2)
\]
are bounded there and have finite one-sided cusp values.
After integration by parts the bulk is the bounded operator
\[
 \mathcal B_{k,q}v(x)
   =\int W(y)(\nu+\partial_y)^q
                         \bigl(k(x-y)/W(y)\bigr)v(y)dy .
 \tag{MA-GW.81}
\]
It is a finite sum of bounded input multipliers followed
by causal convolutions with derivatives of \(k\).
For each derivative order \(j\ge1\), the full boundary
expression on each interval is
\[
 \binom qj\nu^{q-j}
 \sum_{\ell=0}^{j-1}(-1)^{j+\ell}
   [\,\partial_y^\ell(k(x-y)a(y))\,
                          b^{(j-1-\ell)}(y)\,]_{\rm left}^{\rm right}.
 \tag{MA-GW.82}
\]
The endpoint at infinity vanishes.
Keep the lower-endpoint term at zero and the difference
of the two one-sided terms at the cusp.
By (MA-GW.78), \(b\in H^q_{\rm loc}\) and has common
derivative traces through order \(q-1\) there.
Each resulting term is a compact smooth output profile,
a linear combination of \(k^{(i)}(x)\) or \(k^{(i)}(x-L)\),
times a graph-bounded functional \(b^{(j)}(0)\) or
\(b^{(j)}(L)\), \(j<q\).
The one-dimensional Sobolev trace inequality supplies
this bound.  In particular the endpoint and cusp have
not been omitted.

Apply (MA-GW.80)--(MA-GW.82) to \(v=\mathcal R_qf\).
The result is
\(C_k^+P_{U_cZ_{\eta,c}}=\mathcal B_{k,q}\mathcal R_q\)
plus finitely many bounded rank-one maps.
This proves actual ambient trace class.
Conjugating by \(U_c\) and extending by zero proves
the second assertion of (MA-GW.49).
Trace cyclicity with the actual orthogonal projection
identifies this ambient trace with the quotient trace,
already established.  This proves all of (MA-GW.50).

\emph{6. The two orders of actual cutoff limits.}
The exact orthogonal decomposition in \(H\) is
\[
                 Q_{\eta,\Lambda}=P_c-P_{\eta,c}^Z,
                   \qquad P_c=\mathbf1_{[c,\infty)} .
\]
For a compact smooth position multiplier, \(M_\chi C_gP_c\)
is trace class with trace zero.  To verify this before
taking the diagonal, extend the compact-position
convolution kernel to a smooth compact rectangle.
Its rapidly convergent double Fourier expansion is nuclear.
The weight on each compact coordinate projection is bounded
and boundedly invertible, so it only inserts bounded
multipliers.  Restricting by \(P_c\) preserves trace class.
Bounded smooth approximation of that indicator and
trace-norm continuity identify the trace with
\(g(0)\int\chi=0\).
Consequently
\[
 \operatorname{Tr}(M_{\chi_R}C_gQ_{\eta,\Lambda})
     =-\operatorname{Tr}(M_{\chi_R}C_gP_{\eta,c}^Z).
 \tag{MA-GW.83}
\]
For fixed \(\Lambda\), strong convergence of the uniformly
bounded exhaustion paired with (MA-GW.49) gives its limit
in trace norm and hence the first trace in (MA-GW.51).
That trace is already independent of \(\Lambda\).

For the opposite order, the actual all-weight condensation
(MA-CGP.14) gives trace-norm convergence at each fixed
position cutoff.  Its proof uses smooth frequency cutoffs
in the actual character source, both automatic polar
conditions, and strong monotone convergence of the actual
orthogonal projections; it does not use the \(s=0\)
Hardy identification.
The completed global cutoff theorem MA-WP-003 then gives
the actual full-strip trace, with interior multiplicity
\(m_\eta(\rho)\) and exactly \(2j<s-1\) at each edge.
These are precisely (MA-GW.52).
This proves the second order in (MA-GW.51).
Subtracting the two absolutely convergent series yields
(MA-GW.53): the upper-edge truncations cancel and only
the omitted lower-edge lengths remain.
\end{proof}

\paragraph{Exact analytic inputs and scope.}
The source formula is (MA-CGP.13), with its preceding
actual finite-vector and Poisson proofs (MA-CGP.5)--(MA-CGP.7).
Those proofs retain the raw compact type, Gauss phase,
real parity, complex additive measure \(2\,dx\,dy\),
and exact conductor scale.
For Dirichlet characters \(Q=q\); for the non-induced
cubic character \(Q=3\cdot269^2\).
The constant in (MA-GW.59) is therefore fully specified,
rather than absorbed into a changed character divisor.
No archimedean moment or canceled-pole factor belongs
to the present nontrivial character range.

The unconditional fixed-character polynomial growth,
local count and Euler-center bound are NT-Z10,
(NT-J-upper)--(NT-J-lower), in the supplied number-theory
module.  The complete good-height and disk minimum-modulus
proof is MT-M18, (M18.5)--(M18.8).
These arguments were reread for the present extension;
they concern a fixed Hecke function and give no uniformity
in conductor or field.
The inherited primary completed-function locators are
Ramakrishnan--Valenza, \emph{Fourier Analysis on Number Fields},
\S7.4, Theorem~7-19, printed p.~279, and \S7.7,
Theorem~7-28, printed p.~289.
No GRH statement is used to obtain those analytic bounds.

The finite-length modules here are the proved finite-rank
principal parts of an actual weighted quotient resolvent.
Their trace-class smoothing, graph-domain lift and cutoff
limits precede evaluation of the zero sums.
For the character's rank-one finite-monodromy Artin data,
and the non-induced cubic representation, the exact
Grothendieck/Artin placement remains MA-CH and MA-CQ.
For the inherited trace framework the available locators
remain cote~1 scans~079--080 and~111--112,
ALPOGE\_BRIDGES BR-TR-01/02.
No new archive inspection, arithmetic positivity,
half-integral Tate object, or zero-location theorem is
asserted by this extension.

All norm comparisons are for a fixed finite band.
No uniform trace-norm estimate as \(\Lambda\to\infty\),
joint-limit interchange, or equality with the separately
constructed local character Fourier cutoff follows from
this theorem.
